\chapter{Inleiding}
% TIP:
% - Introduceer Solvware en SolvCRM.
% - Leg het belang uit van automatisering van boekhoudprocessen.
% - Wat is er nu handmatig, wat kost tijd?
% - Wat wil de organisatie bereiken?
% Eindig met de structuur van dit document.

Solvware is een softwarebedrijf gevestigd in Waalwijk.
Het biedt als hoofdproduct SolvCRM aan, een alles-in-één CRM- en ERP-systeem voor bedrijven.
Het systeem biedt mogelijkheden voor boekhouding, zoals het bijhouden van klantrelaties, facturen en betalingen.
Omdat klanten van Solvware vaak al bestaande boekhoudsoftware gebruiken, wil Solvware het eenvoudig maken om deze gegevens tussen systemen te synchroniseren.\\
Door een koppeling te realiseren tussen SolvCRM en e-Boekhouden.nl kunnen boekhoudgegevens automatisch worden gesynchroniseerd.
Dit bespaart de klant tijd en verkleint de kans op menselijke fouten bij de gegevensinvoer.\\
Het doel van dit project is om een koppeling te ontwikkelen tussen SolvCRM en e-Boekhouden.nl die periodiek klantrelaties, verzonden facturen en ontvangen factuurbetalingen synchroniseert.

\section{Leeswijzer}
% TIP: Vertel per hoofdstuk wat de lezer kan verwachten.
In dit verslag wordt stap voor stap toegewerkt naar het antwoord op de centrale vraag.
Hieronder wordt per hoofdstuk kort toegelicht wat de lezer kan verwachten.\\
De Inleiding wordt vervolgd door de Aanleiding en Probleemstelling.
In dit hoofdstuk wordt de aanleiding van dit onderzoek en de bijbehorende probleemstelling uitgelegd.
Het hoofdstuk hierna behandelt de doelstelling en onderzoeksvragen.
Hier wordt duidelijk beschreven wat het doel is van dit onderzoek, wat de hoofdvraag is, en welke deelvragen hier bij horen.
Bij het hoofdstuk onderzoeksmethoden, wordt per deelvraag beschreven welke onderzoeksmethoden worden toegepast, en hoe eventuele resultaten zijn verwerkt.
Vervolgens worden de resultaten per deelvraag getoond.
Hierna volgen de conclusie, eventuele discussie en aanbevelingen.
Het document eindigt met een reflectie op het stageonderzoek.
